\section{Contextual Bandits and the Bridge to MDPs}

% --- Setup + interaction merged into one slide ---
\begin{frame}{Contextual Bandits}
\textbf{The one change:} before choosing, the agent is shown a \emph{context} $s$ --- a user profile, a time of day, a search query.
\vspace{0.4em}
\begin{itemize}
    \setlength{\itemsep}{0.55em}
    \item Each step: the environment draws a context $s_t$ $\;\to\;$ the agent picks $a_t$ $\;\to\;$ the reward comes from $\mathcal{R}^{a_t}_{s_t}$, which now depends on \emph{both}
    \item The best arm is no longer a single answer --- it \textbf{depends on the context}. Every quantity from today gains an $s$: $Q(a) \to Q(s,a)$
    \item \textbf{Is the context a state?} \textcolor{red}{Yes --- but a one-step one.} In an MDP your action changes what you see next, $s_{t+1} \sim P(\cdot \mid s_t, a_t)$. Here the next context is drawn \emph{independently of what you did}: your action moves the reward, never the future
    \item That is exactly why there is no credit assignment and no Bellman equation here --- and why contextual bandits are deployed far more widely than full RL
    \item \textcolor{gray}{News recommendation, ad selection, treatment assignment --- the right answer is almost never the same for everyone.}
\end{itemize}
\end{frame}

% --- LinUCB: 3 slides collapsed to 1, motivated by the sparsity table ---
\begin{frame}{LinUCB: Sharing Data Across Contexts}
\textbf{The problem.} One estimate per (context, action) cell is hopeless --- the table is almost entirely empty.
\begin{center}
\footnotesize
\renewcommand{\arraystretch}{1.15}
\begin{tabular}{lccc}
\hline
\textbf{User} & \textbf{Article A} & \textbf{Article B} & \textbf{Article C} \\
\hline
sports, 25, US   & 4 / 50 & --- & --- \\
sports, 27, UK   & ---    & --- & --- \\
finance, 51, US  & ---    & 1 / 12 & --- \\
\hline
\end{tabular}\\[0.2em]
{\scriptsize\textcolor{gray}{Cells show \emph{clicks / times shown}; ``---'' means never shown.}}
\end{center}
\begin{itemize}
    \setlength{\itemsep}{0.3em}
    \item \textbf{The move:} describe each cell by \emph{features} $x_{s,a}$ and model the mean reward as linear, $\mathbb{E}[r \mid s,a] = x_{s,a}^\top\theta_a$. Rows 1 and 2 are nearly the same feature vector, so row 1's clicks now inform row 2
    \item \textcolor{gray}{Not copying cells across --- we fit \emph{one weight vector per arm} on all the data at once, and read every cell off those weights}
    \item \textbf{Same optimism as UCB1}, but the bonus is large in \emph{feature directions} short of data: we are uncertain about \emph{finance readers}, not about one empty cell
\end{itemize}
\centering\small\textcolor{gray}{Yahoo front page (Li et al., WWW 2010): 33M events, 12.5\% click lift over a context-free baseline.}
\end{frame}

\begin{frame}{Bandits $\longrightarrow$ MDPs}
\begin{center}
\renewcommand{\arraystretch}{1.5}
\large
\begin{tabular}{ll}
One-armed bandit     & $=$ \emph{one-state} MDP \\
Contextual bandit    & $=$ \emph{one-step} MDP \\
Add state transitions $\Longrightarrow$ & the full MDP from Day 1 \\
\end{tabular}
\end{center}
\vspace{0.8em}
\begin{itemize}
    \setlength{\itemsep}{0.5em}
    \item Everything we just covered --- UCB, Thompson, gradient bandits --- \emph{lifts} to MDPs in principle
    \item Naively, the state space explodes: visiting every $(s, a)$ enough times isn't feasible in deep RL
    \item \textcolor{red}{\textbf{Up next:}} how to scale these exploration ideas to deep RL with sparse rewards.
\end{itemize}
\end{frame}
